
\documentclass[11pt]{article}

\usepackage[margin=1in]{geometry}
\usepackage{amsmath,amssymb,mathtools}
\usepackage{booktabs}
\usepackage{enumitem}
\usepackage[hidelinks]{hyperref}
\usepackage[T1]{fontenc}
\usepackage[utf8]{inputenc}

\title{Rotoexpando Geometry: Definitions and Characterization\\
\large Nested Helices, Framed Spheres, Closure, and Equivalent Representations}
\author{Nathan McKnight and collaborators}
\date{11 September 2026}

% Portable manuscript symbol for Hagalaz.
% In a Unicode/XeLaTeX build this may be replaced visually by the rune ᚼ.
\newcommand{\Hag}{\mathsf{H}}
\newcommand{\SO}{\mathrm{SO}}
\newcommand{\Conf}{\mathrm{Conf}}
\newcommand{\R}{\mathbb{R}}
\newcommand{\Sph}{\mathbb{S}}

\begin{document}
\maketitle

\begin{abstract}
We define a geometry of nested rotating offsets whose simplest realization is a
fixed-radius helix and whose higher-order realizations include helices carried by
helices, framed-sphere comparisons, toroidal encodings, wave/phase descriptions,
and braid-like closure structures.  The central object is not assumed to be a
physical theory.  It is a geometric relation that may be represented by points,
line segments, framed spheres, ribbons, screw motions, phase oscillators, or
equivalent normalized constructions.  The main quantities are adjacent-order
scale ratio, centerpoint offset, relative rotation, phase ratio, chirality, and
finite-core thickness.  Special configurations arise at closures, degeneracies,
reversals, tangencies, self-contact thresholds, and similarity-return points.
The notation ``Hagalaz'' is used only as a compact mathematical shorthand for a
declared rotoexpando relation; it is optional and carries no independent
ontological content.
\end{abstract}

\section{Scope and guiding convention}

The construction is studied first as pure geometry.  Standard mathematics may be
used freely whenever its assumptions are stated and satisfied.  Physical
interpretations imported from external theories are not part of the definition.

The default baseline is:
\begin{enumerate}[label=(\roman*)]
    \item smooth evolution;
    \item isotropic and symmetric growth or propagation where applicable;
    \item fixed helical radius within a given order unless explicitly varied;
    \item constant pitch/phase rate in the canonical case;
    \item fixed chirality until a reversal is explicitly introduced;
    \item no distortion, self-crossing, or asymmetric forcing until those are
          introduced as controlled extensions.
\end{enumerate}

A crucial distinction is made between \emph{growth} and \emph{expansion}.
A helix may extend indefinitely by screw propagation while retaining fixed
diameter.  Dynamic radial dilation is therefore optional, not fundamental.
Inter-order scale ratios may be non-unit even when no object is expanding in
time.

\section{Primitive point-pair representation}

Let
\[
P(s),Q(s)\in\R^d
\]
be two points.  Their relative vector is
\[
\mathbf{h}(s)=Q(s)-P(s).
\]
Define
\[
\ell(s)=\|\mathbf{h}(s)\|,
\qquad
\widehat{\mathbf{h}}(s)=\frac{\mathbf{h}(s)}{\|\mathbf{h}(s)\|}.
\]

Thus a two-point relation contains:
\begin{itemize}
    \item a length $\ell$;
    \item a direction $\widehat{\mathbf h}$;
    \item in three dimensions, two directional angles.
\end{itemize}

The oriented line segment is a representation of the same data:
\[
L(s,u)=P(s)+u\,\mathbf h(s),
\qquad 0\le u\le1.
\]

A two-point relation does \emph{not} by itself encode roll about its own axis.
Axial twist requires an additional marker: for example a third point, a ribbon,
a material cross-section, or an explicitly transported transverse frame.

\section{Canonical one-order helix}

For a straight carrier axis $\mathbf e_3$, the canonical circular helix is
\[
\gamma(\theta)
=
\rho\cos\theta\,\mathbf e_1
+
\rho\sin\theta\,\mathbf e_2
+
h\theta\,\mathbf e_3,
\]
where $\rho$ is the fixed radial offset and $h$ is axial advance per radian.
This is the orbit of a one-parameter screw motion: rotation plus translation at
fixed offset.

Its curvature and torsion are
\[
\kappa=\frac{\rho}{\rho^2+h^2},
\qquad
\tau=\frac{h}{\rho^2+h^2},
\]
so
\[
\frac{\tau}{\kappa}=\frac{h}{\rho}.
\]

If $\beta$ is the pitch angle measured relative to the plane orthogonal to the
carrier axis, then
\[
\tan\beta=\frac{h}{\rho}
=\frac{\tau}{\kappa}.
\]

The degenerate limits are:
\[
h=0 \;\Rightarrow\; \text{circle},
\qquad
\rho=0 \;\Rightarrow\; \text{straight line}.
\]
A sign change of $h$ or of angular progression reverses handedness.

\section{Nested orders}

Let order $n$ be carried by order $n+1$.  Locally, the tangent or axis of the
higher-order carrier supplies the axis about which the lower-order structure
winds.  One may therefore define recursively
\[
A_{n+1}
\longrightarrow
T_n
\longrightarrow
N_n
\longrightarrow
T_{n-1}
\longrightarrow\cdots,
\]
where $T_n$ is the local tangent of order $n$ and $N_n$ denotes the normal-plane
data associated with that tangent.

For adjacent orders, retain the dimensionless quantities
\[
\mu_n=\frac{D_{n+1}}{D_n},
\qquad
\nu_n=\frac{k_{n+1}}{k_n},
\qquad
\xi_n=\frac{\|c_{n+1}-c_n\|}{D_n},
\]
together with relative phase $\Delta\phi_n$, relative frame rotation $Q_n$, and
chirality data $\chi_n$.

Here:
\begin{itemize}
    \item $\mu_n$ is an inter-order diameter ratio;
    \item $\nu_n$ is an inter-order wavenumber/phase-rate ratio;
    \item $\xi_n$ is normalized centerpoint offset.
\end{itemize}

No equality between $\mu_n$ and $\nu_n$ is assumed.

\section{Framed spheres and the Universal Indicatrix representation}

A helical order may generate a sphere by rotating its local radial segment
through all allowed orientations about the relevant center.  Represent a framed
sphere by
\[
\mathcal S_i=(c_i,r_i,F_i),
\qquad
F_i\in\SO(d).
\]

For two framed spheres, the relative similarity data are
\[
\sigma_{ij}=\frac{r_j}{r_i},
\qquad
Q_{ij}=F_jF_i^{-1},
\qquad
\Delta c_{ij}=c_j-c_i.
\]

After centering and normalizing the radii, both spheres may be represented on the
same unit sphere while retaining two generally misaligned frames.  The
``cockeyed'' superposed grids therefore carry the relative angular information
that a bare sphere, being rotationally symmetric, cannot retain.

This normalized framed-sphere comparison is the core Universal Indicatrix (UI)
representation used here.

The previously studied Three-Spheres construction is treated as a special
constrained configuration inside the more general UI state space.

\section{Hagalaz / rotoexpando notation}

``Hagalaz'' is optional shorthand.  In this manuscript it is rendered
portably as $\Hag$.

For adjacent normalized orders,
\[
\Hag_{ij}:=(\sigma_{ij},Q_{ij}),
\]
or, if translation is retained,
\[
\Hag_{ij}:=(\Delta c_{ij},\sigma_{ij},Q_{ij}).
\]

The corresponding Euclidean similarity map is
\[
x\mapsto c_j+\sigma_{ij}Q_{ij}(x-c_i).
\]

Within a single fixed-diameter helical order, the canonical evolution need not
contain dilation at all: it may be a screw motion with fixed offset.  The scale
factor in $\Hag_{ij}$ is therefore primarily an \emph{inter-order comparison}
unless dynamic dilation is explicitly introduced.

\section{Equivalent representations}

The same relation may be encoded in several mathematically equivalent or
losslessly convertible forms, provided the required frame data are retained:

\[
\text{point pair}
\;\leftrightarrow\;
\text{oriented segment}
\;\leftrightarrow\;
\text{screw/helix}
\;\leftrightarrow\;
\text{framed sphere/UI}
\;\leftrightarrow\;
\text{phase oscillator}
\;\leftrightarrow\;
\text{toroidal winding}.
\]

A ribbon or framed segment contains one additional roll/twist degree of freedom
that a bare two-point segment does not.

For a planar complex representation,
\[
x+iy=\rho e^{i\theta}.
\]
For a traveling phase,
\[
x+iy=\rho e^{i(ks-\omega t+\phi)}.
\]

If the same propagation speed is used for compared modes, then
\[
\omega=vk,
\qquad
k=\frac{2\pi}{\lambda},
\]
and therefore
\[
\frac{\lambda_m}{\lambda_n}
=
\frac{k_n}{k_m}
=
\frac{\omega_n}{\omega_m}.
\]

Thus phase-rate ratios, wavenumber ratios, frequency ratios, and wavelength
ratios are interchangeable descriptions under the declared common-speed
assumption.

\section{Toroidal encoding and Whirlygig reduction}

A helix compressed into its transverse plane loses its axial coordinate only if
that coordinate is discarded.  It may instead be stored as a second periodic
phase.

Let
\[
\phi=\phi(s),
\qquad
\psi=\psi(s),
\]
and define
\[
X(\phi,\psi)
=
\bigl(
(R+a\cos\psi)\cos\phi,\,
(R+a\cos\psi)\sin\phi,\,
a\sin\psi
\bigr).
\]

Then a one-parameter helical history may be encoded as a curve on a torus:
\[
s\mapsto X(\phi(s),\psi(s)).
\]

The ratio
\[
q=\frac{d\psi}{d\phi}
\]
is the toroidal winding slope.  Rational $q$ yields finite closure; irrational
$q$ yields no exact finite closure in the ideal mathematical limit.

The Whirlygig is treated as a special lower-dimensional
compression/projection of this broader phase geometry.  Its exact status is to
be established by explicit transformation rather than assumed.

\section{Closure and ``click points''}

For order $n$, define a closure event by
\[
\theta_n=2\pi j.
\]
At those events, an adjacent order with phase ratio $\nu$ has phase
\[
\Theta_j
=
\Theta_0+2\pi j\nu
\pmod{2\pi}.
\]

If
\[
\nu=\frac{p}{q}
\]
in lowest terms, then there are $q$ distinct closure states before exact
recurrence.  Irrational $\nu$ produces no exact finite recurrence.

It is useful to distinguish:
\begin{itemize}
    \item \textbf{Euclidean closure}: the configuration returns to the same
          position, orientation, and scale;
    \item \textbf{similarity closure}: the normalized shape/orientation returns
          while the absolute scale may differ;
    \item \textbf{partial-phase closure}: one phase closes while another does not;
    \item \textbf{near closure}: phases return within a specified tolerance.
\end{itemize}

The term ``holonomy'' should be reserved for cases where a frame or state is
transported around a genuinely closed loop.  Otherwise ``return map'',
``recurrence'', ``monodromy'', or ``similarity closure'' may be more precise.

\section{Isotropic centerpoint-offset layouts}

At a closure state with angular coordinate $\Theta_j$, place a segment base at
\[
b_j=d\,e^{i\Theta_j},
\]
where $d$ is the centerpoint offset magnitude.  Let the segment have length
$\ell$ and relative angle $\gamma_j$.  Its endpoint is
\[
z_j
=
e^{i\Theta_j}
\left(
d+\ell e^{i\gamma_j}
\right).
\]

For constant $d,\ell,\gamma$,
\[
|z_j|
=
\sqrt{d^2+\ell^2+2d\ell\cos\gamma},
\]
so all endpoints lie on a concentric circle.

Special cases include
\[
\gamma=0
\quad\Rightarrow\quad
|z|=d+\ell,
\]
\[
\gamma=\pi
\quad\Rightarrow\quad
|z|=|d-\ell|,
\]
and the complete cancellation condition
\[
d=\ell,\qquad\gamma=\pi
\quad\Rightarrow\quad
z=0.
\]

A radial crossover occurs when
\[
d+\ell\cos\gamma=0.
\]

If the relative angle itself rotates,
\[
\gamma_j=\gamma_0+m\Theta_j,
\]
then
\[
z_j
=
d\,e^{i\Theta_j}
+
\ell e^{i((m+1)\Theta_j+\gamma_0)},
\]
a two-frequency phasor sum that generates rosette/trochoidal families and
provides a natural route toward Whirlygig-like planar traces.

\section{Recursive inter-order layouts}

For a constant adjacent-order scale ratio $\mu$ and angular offset $\alpha$, a
recursive segment chain may be written
\[
z_N
=
z_0+
D_0\sum_{j=0}^{N-1}\mu^j e^{ij\alpha}.
\]
Hence
\[
z_N
=
z_0+
D_0
\frac{1-(\mu e^{i\alpha})^N}
     {1-\mu e^{i\alpha}},
\]
provided $\mu e^{i\alpha}\neq1$.

This immediately distinguishes:
\[
0<\mu<1
\quad\Rightarrow\quad
\text{contracting hierarchy},
\]
\[
\mu=1
\quad\Rightarrow\quad
\text{scale-preserving hierarchy},
\]
\[
\mu>1
\quad\Rightarrow\quad
\text{expanding hierarchy}.
\]

If
\[
\alpha=\frac{2\pi p}{q},
\qquad
\mu=1,
\]
then after $q$ steps
\[
z_q=z_0.
\]

If the same angular closure occurs with $\mu\neq1$, orientation may recur while
scale changes by $\mu^q$, giving a similarity-return point.

\section{Dimensional measure scaling}

If adjacent linear scales differ by $\mu$, then every $k$-dimensional Euclidean
measure scales as
\[
M^{(k)}_{n+1}=\mu^k M^{(k)}_n.
\]

Thus in three dimensions:
\[
\frac{D_{n+1}}{D_n}=\mu,
\qquad
\frac{A_{n+1}}{A_n}=\mu^2,
\qquad
\frac{V_{n+1}}{V_n}=\mu^3.
\]

For $S^3\subset\R^4$, the hypersurface measure scales as $\mu^3$, while enclosed
four-volume scales as $\mu^4$.

The two-vector relation
\[
R_*^2=d^2+\ell^2+2d\ell\cos\gamma
\]
is dimension-independent in Euclidean space.

\section{Four-dimensional lift}

A generic rotation in $\R^4$ may be decomposed, in a suitable orthogonal basis,
into rotations in two mutually orthogonal planes:
\[
Q
=
\begin{pmatrix}
R(\theta_1)&0\\
0&R(\theta_2)
\end{pmatrix}.
\]

Equivalently,
\[
z_1=A_1e^{i\theta_1},
\qquad
z_2=A_2e^{i\theta_2}.
\]

Complete angular recurrence requires simultaneous closure:
\[
m\theta_1=2\pi p,
\qquad
m\theta_2=2\pi q
\]
for integers $m,p,q$.

This creates exact notions of:
\begin{itemize}
    \item full closure;
    \item partial closure in one rotation plane;
    \item near-simultaneous closure;
    \item irrational/quasiperiodic nonclosure.
\end{itemize}

\section{Finite thickness and first-contact constraints}

For a finite tube of radius $a$ around a centerline $X(s)$, the first
nonlocal-contact condition is
\[
\inf_{s\not\approx t}\|X(s)-X(t)\|=2a.
\]

Local curvature imposes a separate constraint of the schematic form
\[
a<\frac{1}{\kappa_{\max}}
\]
for the simplest embedded-tube setting.

Thus the admissible finite-core radius is controlled jointly by local curvature
and nonlocal self-approach.  The regular kinematic regime terminates at a purely
geometric threshold before any model-specific collision, deformation,
reconnection, penetration, or exclusion rule is supplied.

For the circular helix,
\[
D^2(\Delta)
=
2\rho^2(1-\cos\Delta)+h^2\Delta^2,
\]
so stationary candidates for nonlocal closest approach obey
\[
\rho^2\sin\Delta+h^2\Delta=0.
\]

\section{Braids and spherical braids}

If multiple strands intersect a reference sphere, their intersection points
\[
u_i(s)\in\Sph^2
\]
define a configuration
\[
U(s)=\bigl(u_1(s),\ldots,u_N(s)\bigr)
\in \Conf_N(\Sph^2)
\]
as long as no two points coincide.

A closed history in this configuration space defines a spherical braid class.
The collision set
\[
u_i=u_j
\]
is precisely where the configuration exits the collision-free configuration
space and where braid-changing events may become geometrically possible.

Thus metric deformation may proceed continuously while braid topology remains
fixed up to the first collision/contact threshold.

\section{Special loci in parameter space}

The pure-geometry program is to map the parameter space without privileging any
numerical value in advance.  Of particular interest are loci corresponding to:

\begin{itemize}
    \item loss of apparent helicity;
    \item handedness reversal;
    \item exact and partial closures;
    \item similarity-return points;
    \item axis coincidence or orthogonality;
    \item tangencies and crossovers;
    \item local curvature saturation;
    \item first nonlocal self-contact;
    \item braid-class change opportunities;
    \item recursive exclusion sets;
    \item transitions between finite, quasiperiodic, dense, and self-similar
          closure structures.
\end{itemize}

Golden-ratio, Fibonacci, Cantor-like, or other special structures are therefore
treated as hypotheses to be discovered from the special-locus map, not inserted
as assumptions.

\section{Candidate minimal atlas}

For an adjacent pair of orders, a useful compact state is
\[
\boxed{
\mathcal C_n
=
\left(
\mu_n,\,
\nu_n,\,
\xi_n,\,
\Delta\phi_n,\,
Q_n,\,
\chi_n,\,
a_n/D_n
\right)
}
\]
with:
\[
\mu_n=\frac{D_{n+1}}{D_n},
\qquad
\nu_n=\frac{k_{n+1}}{k_n},
\qquad
\xi_n=\frac{\|c_{n+1}-c_n\|}{D_n}.
\]

This atlas is intended to be convertible among point-pair, segment, helix,
framed-sphere/UI, wave, torus, ribbon, and braid representations.

\section{Immediate mathematical program}

The next stage is deliberately computational and classificatory:

\begin{enumerate}
    \item define the exact nested-helix-to-framed-sphere map;
    \item derive the induced adjacent-order UI transformation;
    \item verify lossless conversion among the equivalent representations;
    \item sweep the minimal parameter set
    \[
    (\mu,\nu,\xi,\gamma,a/D);
    \]
    \item locate special loci without fitting target values;
    \item classify each locus by standard mathematical name where one exists;
    \item identify any residual construction or invariant for which no standard
          named object is found;
    \item only after the geometry is stable, investigate possible H(s)H
          applications.
\end{enumerate}

\section{Nonclaim}

This document defines and characterizes a geometric construction.  It does not
assert that any particular physical system is governed by this geometry, nor
that special numerical values must occur.  External mathematical machinery may
be used wherever applicable; external physical theories are not silently
imported as premises.

\end{document}
